\section{式の synthesis 規則}

以下では $q;S;\Gamma\vdash e\Rightarrow\tau\dashv q';S'$ を
$\DDSynth_{\Sighat}(q,S,\Gamma,e,\tau,q',S')$ の略記とする．
$\alpha_{q_\tau}$ は supply $q$ が指す次の target variable である．

\subsection{変数・関数・適用}

変数 lookup は prevailing substitution を context に適用してから scheme を fresh instantiate
する．lambda domain は任意に選ばず，必ず fresh metavariable とする．関数適用は function を
先に function type へ align し，その後で argument を domain に対して check する．
\begin{mathpar}
\inferrule{
  (S\Gamma)(x)=\sigma \\
  \Inst_q(\sigma)=(\tau,q')
}{
  q;S;\Gamma\vdash x\Rightarrow\tau\dashv q';S
}\rn{DD-VAR}

\inferrule{
  q+\langle0,1\rangle;S;\Gamma,x:\mono{\alpha_{q_\tau}}\vdash
    e\Rightarrow\tau\dashv q';S'
}{
  q;S;\Gamma\vdash\lambda x.e\Rightarrow
    \alpha_{q_\tau}\to\tau\dashv q';S'
}\rn{DD-LAM}

\inferrule{
  q;S;\Gamma\vdash e_1\Rightarrow\tau_f\dashv q_1;S_1 \\
  \DDAlignTypes(S_1,\tau_f,
    \alpha_{(q_1)_\tau}\to\alpha_{(q_1)_\tau+1},S_2) \\
  q_1+\langle0,2\rangle;S_2;\Gamma\vdash e_2\Leftarrow
    \alpha_{(q_1)_\tau}\dashv q_2;S_3
}{
  q;S;\Gamma\vdash e_1\,e_2\Rightarrow
    \alpha_{(q_1)_\tau+1}\dashv q_2;S_3
}\rn{DD-APP}
\end{mathpar}

この順序により，function domain が slot に解決された場合，argument の構文に関係なく
その一つの $\DDCheck$ cut で matcher-to-slot coercion を選べる．

\subsection{リテラル・タプル・コンストラクタ}

\begin{mathpar}
\inferrule{ }{
  q;S;\Gamma\vdash n\Rightarrow\mathsf{Int}\dashv q;S
}\rn{DD-LIT}

\inferrule{
  \DDSynths_{\Sighat}(q,S,\Gamma,\overline e,
    \overline\tau,q',S')
}{
  q;S;\Gamma\vdash(\overline e)\Rightarrow
    (\overline\tau)\dashv q';S'
}\rn{DD-TUPLE}

\inferrule{
  \Inst_q(\Sighat_D(C))=(\overline\rho\to\rho,q_0) \\
  \DDChecks_{\Sighat}(q_0,S,\Gamma,\overline e,
    \overline\rho,q',S')
}{
  q;S;\Gamma\vdash C\,\overline e\Rightarrow
    \rho\dashv q';S'
}\rn{DD-CON}

\inferrule{
  \Inst_q(\Sighat_{\mathit{prim}}(op))=
    (\overline\rho\to\rho,q_0) \\
  \DDChecks_{\Sighat}(q_0,S,\Gamma,\overline e,
    \overline\rho,q',S')
}{
  q;S;\Gamma\vdash op(\overline e)\Rightarrow
    \rho\dashv q';S'
}\rn{DD-PRIM}
\end{mathpar}
$\DDSynths$ と $\DDChecks$ は head の出力 state を tail へ渡すため，constructor と primitive
の引数も左から右へ処理される．

\subsection{Let と再帰}

\begin{mathpar}
\inferrule{
  q;S;\Gamma\vdash e_1\Rightarrow\tau_1\dashv q_1;S_1 \\
  \sigma=\Gen(\Sighat,S_1\Gamma,S_1\tau_1) \\
  q_1;S_1;\Gamma,x:\sigma\vdash e_2\Rightarrow
    \tau_2\dashv q';S'
}{
  q;S;\Gamma\vdash
  \mathtt{let}\ x=e_1\ \mathtt{in}\ e_2\Rightarrow
    \tau_2\dashv q';S'
}\rn{DD-LET}

\inferrule{
  f\ne x \\
  \DirectSelf(f,e)\quad\mathsf{NonMatcherBody}(e) \\
  q+\langle0,2\rangle;S;\Gamma,
  f:\mono{(\alpha_{q_\tau}\to\alpha_{q_\tau+1})},
  x:\mono{\alpha_{q_\tau}}\vdash e\Rightarrow\rho\dashv q_1;S_1 \\
  \DDAlignTypes(S_1,\rho,\alpha_{q_\tau+1},S')
}{
  q;S;\Gamma\vdash\mathtt{fix}\ f\ x.e\Rightarrow
    \alpha_{q_\tau}\to\alpha_{q_\tau+1}\dashv q_1;S'
}\rn{DD-FIX}
\end{mathpar}
$\DirectSelf(f,e)$ は shadowing-aware な構文条件である．$f$ の自由出現を application spine の
head に限り，bare value，引数位置，alias 経由，高階の受渡しを拒否する．matcher literal を
body に持つ recursion は第5章の専用 placeholder 規則を使う．

\subsection{標準 matcher producer}

\begin{mathpar}
\inferrule{ }{
  q;S;\Gamma\vdash\mathtt{something}\Rightarrow
    \Matcher{\Any}{\alpha_{q_\tau}}\dashv
    q+\langle0,1\rangle;S
}\rn{DD-SOMETHING}
\end{mathpar}
$\mathtt{something}$ 自身は matcher producer である．slot が必要かどうかはこの規則では
決めず，外側の $\DDCheck$ cut が expected head を見て決める．
